% =========================================================================
% CAPÍTULO 4
% =========================================================================
\chapter{Arquitectura Computacional del Sistema}

La Figura \ref{fig:arquitectura} sintetiza el flujo completo de datos, desde las fuentes regulatorias hasta la generación de alertas de riesgo.

\begin{figure}[H]
\centering
\begin{tikzpicture}[
    node distance=0.55cm and 1.4cm,
    box/.style={draw, rounded corners, fill=boxbg, minimum width=3.6cm, minimum height=0.9cm, align=center, font=\small},
    arr/.style={-{Stealth}, thick}
]
\node[box] (fuentes) {Fuentes de datos\\(CNBV, Banxico)};
\node[box, below=of fuentes] (ingesta) {Ingesta\\(Excel/Parquet)};
\node[box, below=of ingesta] (limpieza) {Limpieza e imputación\\(DuckDB / SQL vectorizado)};
\node[box, below=of limpieza] (features) {Ingeniería de características\\(Polars: rezagos, medias móviles, ratios)};
\node[box, below=of features] (dataset) {Dataset panel final\\(1189 obs. $\times$ 46 instituciones)};
\node[box, below=of dataset] (modelos) {Modelos\\(ARIMA/SARIMA, XGBoost, LightGBM, RF, Stacking)};
\node[box, below=of modelos] (validacion) {Validación temporal\\(TimeSeriesSplit, partición por posición)};
\node[box, below=of validacion] (shap) {Interpretabilidad\\(SHAP)};
\node[box, below=of shap] (alerta) {Sistema de alerta temprana\\(score de riesgo)};
\node[box, below=of alerta] (resultados) {Resultados y reporte\\(figuras, tablas, dashboard)};

\draw[arr] (fuentes) -- (ingesta);
\draw[arr] (ingesta) -- (limpieza);
\draw[arr] (limpieza) -- (features);
\draw[arr] (features) -- (dataset);
\draw[arr] (dataset) -- (modelos);
\draw[arr] (modelos) -- (validacion);
\draw[arr] (validacion) -- (shap);
\draw[arr] (shap) -- (alerta);
\draw[arr] (alerta) -- (resultados);
\end{tikzpicture}
\caption{Arquitectura general del pipeline de alerta temprana. Fuente: elaboración propia.}
\label{fig:arquitectura}
\end{figure}

El procesamiento tradicional basado en estructuras densas sin tipado estricto genera problemas de memoria y lentitud en agregaciones de panel. La arquitectura implementada emplea DuckDB como motor analítico embebido que ejecuta SQL vectorizado directamente sobre archivos Parquet sin necesidad de cargar la totalidad de los datos en memoria, y Polars —basado en Apache Arrow y escrito en Rust— para las transformaciones de ingeniería de características, aprovechando paralelismo multihilo sin las limitaciones del \textit{Global Interpreter Lock} de Python.

En cuanto a la reproducibilidad, todos los modelos estocásticos fijan semilla aleatoria (\texttt{random\_state=42}). El código completo, las versiones de librerías utilizadas y las instrucciones de ejecución se documentan en el Anexo C. El pipeline completo —desde la lectura del archivo Parquet hasta la generación de las 17 figuras y las métricas reportadas en el Capítulo 6— es ejecutable de manera autocontenida y determinista.

